\subsection{The Row-by-Column Rule}

\begin{definitionbox}[Product of two matrices]
Let $A=(a_{ij})$ be an $m\times n$ matrix and let $B=(b_{ij})$ be an
$n\times p$ matrix. Their product is the $m\times p$ matrix $C=AB$ whose
entries are
\[
c_{ij}=a_{i1}b_{1j}+a_{i2}b_{2j}+\cdots+a_{in}b_{nj}.
\]
Thus $c_{ij}$ is obtained from row $i$ of $A$ and column $j$ of $B$.
\end{definitionbox}

\begin{interpretationbox}[Why the inner dimensions must match]
A row of $A$ contains $n$ entries. It can be paired with a column of $B$ only
when that column also contains $n$ entries. Therefore
\[
(m\times n)(n\times p)=m\times p.
\]
The inner dimensions test whether multiplication is possible; the outer
dimensions determine the size of the result.
\end{interpretationbox}

\begin{examplebox}[Computing a product]
Let
\[
A=\begin{pmatrix}1&3\\-2&4\end{pmatrix},
\qquad
B=\begin{pmatrix}2&0\\5&-1\end{pmatrix}.
\]
The first entry of $AB$ is
\[
1\cdot2+3\cdot5=17,
\]
while the first row and second column give
\[
1\cdot0+3\cdot(-1)=-3.
\]
Continuing row by column,
\[
AB=
\begin{pmatrix}
17&-3\\
16&-4
\end{pmatrix}.
\]
\end{examplebox}

\begin{summarybox}[What each product entry records]
Every entry of $AB$ measures one interaction between a row of $A$ and a column
of $B$. Matrix multiplication is therefore a structured combination of whole
rows and columns, not a position-by-position product.
\end{summarybox}
